\documentclass[twoside,openany]{book}
\usepackage[czech]{babel}			% Language
\usepackage[utf8]{inputenc}			% Encoding of characters in this .tex file
\usepackage{cmap}					% Make PDF file searchable and copyable (ASCII characters)
\usepackage{lmodern}				% Make PDF file searchable and copyable (Accented characters)
\usepackage[T1]{fontenc}			% Hyphenate accented words
\usepackage[a6paper, top=0.75cm, bottom=1.5cm, inner=0.75cm, outer=0.75cm]{geometry}		% Paper size and margins
\usepackage[fontsize=9]{scrextend}  % Font size
\usepackage{marginnote}				% Enable margin notes
\usepackage{rotating}				% Rotated environments (for margin headers)
\usepackage[protrusion]{microtype}	% Better typeset results
\usepackage{enumitem}				% Nicer enumeration lists
\usepackage{ragged2e}				% Add \justify command
\usepackage{titlesec}				% Customize chapters
\usepackage[usenames,dvipsnames,svgnames,table]{xcolor}		% Custom colors
\usepackage{tikzsymbols}			% Smileys
\usepackage{chronology}				% Timeline
%\usepackage{fontawesome}			% Special symbols

% ======= TODO ======= 

% Ondras: Napisat uvod k rubrike
% Nedavat vlastnu otazku do vacsiny tyzdnov. Doplnit vzdy strukturovanu otazku (pripadne sede namety na protilahlu stranu) zhruba podla toho, v akom poradi mame naplanovane cvika.
% posunut margin notes "pred/po cviceni", aby boli rovnako daleko od deliacej ciary
% naformatovat "uzitecne nastroje"

% ======= CUSTOM DOCUMENT SETTINGS ========

% Change \chapter and \section command display and spacings
% Source: https://tex.stackexchange.com/questions/63390/how-to-decrease-spacing-before-chapter-title
\titleformat{\chapter}[display]{\normalfont\Large\bfseries}{Týden \thechapter}{0em}{\Large}
\titlespacing*{\chapter}{0em}{-2em}{0em}
\titleformat{\section}[display]{\normalfont\bfseries}{\sectionname}{}{}
\titlespacing*{\section}{0em}{0.5em}{0em}

% \punct command allows to shift the footnotemark above a punctuation
\newlength{\spc} % declare a variable to save spacing value
\newcommand{\punct}[1]{%
  \settowidth{\spc}{#1}% set value of \spc variable to the width of #1 argument
  \addtolength{\spc}{-1.8\spc}% subtract from \spc about two (1.8) of its values making its magnitude negative
  #1% print the optional argument
  \hspace*{\spc}% print an additional negative spacing stored in \spc after #1
}

% Remove headers and position page number in the bottom mid
\pagestyle{plain}

% Smaller spaces between enumeration lists
\setitemize{itemsep=0em}

% Disable automatic indentation, enable paragraph skip
\setlength{\parindent}{0em}
\setlength{\parskip}{0.5em}

% Rule length including the margin
\newlength{\rulelength}
\setlength{\rulelength}{\dimexpr(\textwidth+0.5cm)\relax}

% New table column types with flexible width
%\newcolumntype{C}{>{\centering\arraybackslash}X}
\newcolumntype{C}[1]{>{\centering\arraybackslash}m{#1}}

% Rotated table cell
\newcommand{\rotatedCell}[2][l]{\rotatebox{90}{\renewcommand{\arraystretch}{0.75}\begin{tabular}[#1]{@{}l}#2\end{tabular}}}

% ======= CUSTOM ENVIRONMENTS ========

% Mala poznamka
\newcommand{\note}[1]{\textcolor{gray}{\small\itshape #1}}

% Poznamka za okrajom
\newcommand{\marginsection}[3]{\marginpar{\raisebox{#2\height}{
\begin{turn}{#1}
%\fontfamily{phv}\selectfont
\bfseries \color{gray} \large
#3\end{turn}}}}

% Casova os pre rozvrhnutie planu cvika 
\newcommand{\timeline}{
\vspace*{1.1cm}
\color{gray}
\begin{chronology}[15]{0}{120}{\textwidth}
\end{chronology}
\color{black}
\vspace*{0.9cm}
}

% Smajlici pre sebahodnotenie pocitu z cvika
\newcommand{\smileys}{
Jak jsem spokojen(a) s hodinou?
\raisebox{-0.3em}{\,\,\Changey[2]{-2}\,\,\Changey[2]{-1}\,\,\Changey[2]{0}\,\,\Changey[2]{1}\,\,\Changey[2]{2}}
}

% Zhodnotenie 2 kladov a 2 zaporov cvika
\newcommand{\selfassessment}{
Co se mi podařilo?
\begin{enumerate}
\item
\item
\end{enumerate}

Co mohlo být líp?
\begin{enumerate}
\item
\item
\end{enumerate}
}

% Klasicky tyzden
\newcommand{\commonweek}{
\chapter{}
\vspace*{-2em}

\marginsection{90}{-1}{před výukou}

Příprava:\hspace{1cm} hodin(y)

Jaký je cíl této hodiny? Proč chci studenty učit tohle?
\vspace*{1cm}

Jak poznám, že se cíle podařilo dosáhnout?

\hspace*{-1cm}
\rule{\rulelength}{0.4pt}

\marginsection{90}{-1.2}{po výuce}

\smileys

\selfassessment

Vlastní otázka:

\newpage
\marginsection{-90}{1}{volné poznámky}
}

% Stranka rubriky
\newcommand{\rubricpage}[8]{
\newpage
\begin{tabular}{@{}C{0.48\textwidth}C{0.48\textwidth}@{}l@{}}
\normalsize \bfseries #1 & \normalsize \bfseries #5 & \\[1em] \hline \\[-1.2em]
	\justify \note{Not aware:} #2 & \justify \note{Not aware:} #6 & \\[2em]
$\bullet$ & $\bullet$ & \\
$\bullet$ & $\bullet$ & \\
\justify \note{Novice skill:} #3 & \justify \note{Novice skill:} #7 & \\[3em]
$\bullet$ & $\bullet$ & \\
$\bullet$ & $\bullet$ & \\
$\bullet$ & $\bullet$ & \\
$\bullet$ & $\bullet$ & \\
$\bullet$ & $\bullet$ & \\
\justify \note{Mastery:} #4 & \justify \note{Mastery:} #8 &
\end{tabular}
}

% ======= TITLE PAGE =======

\title{Reflektivní deník učitele}
\author{Kolektiv vyučujících Praktika vedení cvičení}
\date{\today}

% ======= DOCUMENT START =======

\begin{document}
\begin{titlepage}
	\centering
    \vspace*{1cm}
	\includegraphics[width=0.5\textwidth]{img/diary}\par
	\vspace{0.5cm}
	{\huge\bfseries Reflektivní\\ deník učitele\par}
	\vfill
	\raggedright
    \Large
	Majitel:\par
    \vspace{0.2cm}
    Předmět:\par
\end{titlepage}

% ======= INTRODUCTION =======

\chapter*{Jak se z člověka stane dobrý učitel?}

\vspace*{1em}
Položme si namísto statické otázky dynamickou:

\vspace*{1em}
\textit{\large \uv{Co můžeš dělat pro to, aby ses jako učitel zlepšoval?}}

\vspace*{1em}
Ukazuje se, že jsou nutné a nejspíš i dostačující tři věci\punct{.}\footnotemark
\footnotetext{Jsou i další věci, které mohou pomoci (například expertní zpětná vazba) ale ty nejsou dobře škálovatelné.}
\begin{enumerate}
\item \textbf{Pravidelně učit.} \\Potřebuješ mít vlastní zkušenost, nejlépe často a pravidelně opakovanou.
\item \textbf{Reflektovat svoje učení.} \\Všímej si a přemýšlej nad tím, co fungovalo dobře a co příště zkusíš jinak.
\item \textbf{Pozorovat, jak učí ostatní.} \\Přemýšlej o tom, jak to dělají, co jim funguje, co ne a co z toho můžeš použít. 
\end{enumerate}

Tento deníček ti pomůže reflektovat způsob, jakým učíš. Také v~něm nalezneš informace o tom, čeho je užitečné si všímat, jaké otázky si klást a nad čím přemýšlet. Není to však učebnice dobrého učitele -- je to pouze průvodce na cestě, po které už jdeš.

Vhodným doplňkem k vlastní reflexi výuky jsou diskuse s ostatními učiteli o tom, jak učí. Pokud se chceš setkávat s dalšími cvičícími pravidelně a navíc se dozvědět více o tom, co funguje a co ne, zapiš si předmět \textit{DUCIT Praktikum vedení cvičení} na Fakultě informatiky MU.

\newpage
\section*{Jak reflektivní deníček používat?}

Nejlépe pravidelně. 

Při přípravě i po každé své odvedené výuce si zapiš svoje myšlenky. Na třináct týdnů v semestru máš nachystaných třináct dvoustránek a na každé pár nabídnutých otázek. Deníček je malý mimo jiné i proto, aby ti psaní nezabralo moc času.

Za stránkami pro jednotlivé týdny najdeš ještě rubriku učitelských schopností, seznam indikátorů, které je možné na hodině sledovat, připomenutí některých užitečných nástrojů k výuce a na závěr prostor na své vlastní poznámky.

\section*{Jaké jsou cíle používání deníku?}

Chtěli bychom, aby práce s deníkem:
\begin{enumerate}[topsep=0pt]
\item Připomínala učiteli zamyslet se nad svou hodinou (reflexe).
\item Vytvářela prostor pro snadné psaní si poznámek pro příště.
\item Vedla lidi k uvědomění si, co všechno učitel dělá.
\item Umožňovala sledovat svůj vlastní pokrok.
\end{enumerate}

\section*{Máš nápady na zlepšení?}

Budeme rádi za jakékoliv nápady, zkušenosti či poznámky k deníku. Ozvi se prosím emailem Martinovi na \textit{mukrop@mail.muni.cz}.

\newgeometry{top=0.75cm, bottom=1.5cm, inner=0.75cm, outer=1.2cm,
			heightrounded, marginparwidth=0.5cm, marginparsep=0.3cm}

% ======= WEEK 1 =======

\chapter{}
\vspace*{-2em}

\marginsection{90}{-2.1}{před výukou}

Příprava:\hspace{1cm} hodin(y)\\
\note{Příprava materiálů, úkolů, pokynů, ...}

Z jakých bloků se má hodina skládá?\\
\note{Rozvrhni si 2--6 bloků na časovou osu níže.}

\timeline

Jaké precedenty chci ve výuce nastavit? \\
\note{Oslovování, (ne)formálnost, kladení otázek, začátek včas, ...}
\vspace*{1cm}

\hspace*{-1cm}
\rule{\rulelength}{0.4pt}

\marginsection{90}{-1.2}{po výuce}

\smileys
\note{Jaký mám pocit já? Co vidím na studentech?}

\selfassessment

\newpage
\marginsection{-90}{1}{volné poznámky}
\vspace*{-2em}
\note{Jaké chtěné/nechtěné precedenty se objevili na hodině?\\
Jak na další hodině udržím ty chtěné a napravím ty nechtěné?\\
Jaká byla atmosféra během vyučování?\\
Rozumí studenti organizaci kurzu a ví, co se od nich očekává?\\
(Další tipy na otázky najdeš na straně \pageref{indikatory}.)}

% ======= WEEK 2 =======

\chapter{}
\vspace*{-2em}

\marginsection{90}{-2.1}{před výukou}

Příprava:\hspace{1cm} hodin(y)\\
\note{Příprava materiálů, úkolů, pokynů, ...}

Z jakých bloků se má hodina skládá?\\
\note{Rozvrhni si 2--6 bloků na časovou osu níže.}

\timeline

Jaké precedenty chci ve výuce udržet? \\
\note{Oslovování, (ne)formálnost, kladení otázek, začátek včas, ...}
\vspace*{1cm}

\hspace*{-1cm}
\rule{\rulelength}{0.4pt}

\marginsection{90}{-1.2}{po výuce}

\smileys
\note{Jaký mám pocit já? Co vidím na studentech?}

\selfassessment

\newpage
\marginsection{-90}{1}{volné poznámky}
\vspace*{-2em}
\note{Na jaký aspekt výuky se zaměřím při nejbližší hospitaci?\\
Jaký aspekt výuky má sledovat kolega, který navštíví mou hodinu?\\
(Další tipy na otázky najdeš na straně \pageref{indikatory}.)}

% ======= WEEK 3 =======

\chapter{}
\vspace*{-2em}

\marginsection{90}{-1}{před výukou}

Příprava:\hspace{1cm} hodin(y)

Z jakých bloků se má hodina skládá?

\timeline

Jaké otázky na hodině položím studentům?
\vspace*{0.8cm}

\hspace*{-1cm}
\rule{\rulelength}{0.4pt}

\marginsection{90}{-1.2}{po výuce}

\smileys

\selfassessment

Na které otázky studenti reagovali? Na které ne?

\newpage
\marginsection{-90}{1}{volné poznámky}
\vspace*{-2em}
\note{Kolik otázek do publika jsem položil(a)?\\
Kolik otázek položili studenti mně? Je to málo/dost/příliš?\\
(Další tipy na otázky najdeš na straně \pageref{indikatory}.)}

% ======= WEEK 4 =======

\chapter{}
\vspace*{-2em}

\marginsection{90}{-1}{před výukou}

Příprava:\hspace{1cm} hodin(y)

Z jakých bloků se má hodina skládá?

\timeline

\hspace*{-1cm}
\rule{\rulelength}{0.4pt}

\marginsection{90}{-1.2}{po výuce}

\smileys

\selfassessment

\note{Sem napiš otázku, na kterou si chceš po výuce odpovědět.}\\
Vlastní otázka:

\newpage
\marginsection{-90}{1}{volné poznámky}
\vspace*{-2em}
\note{Jaké znalosti a dovednosti chci naučit? (Popiš je do rubriky!)\\
Jaký typ učitele chci být? Co bych měl(a) pro to dělat?\\
(Další tipy na otázky najdeš na straně \pageref{indikatory}.)}

% ======= WEEK 5 =======

\chapter{}
\vspace*{-2em}

\marginsection{90}{-1}{před výukou}

Příprava:\hspace{1cm} hodin(y)

Jaký je cíl této hodiny? Proč chci studenty učit tohle? \\
\note{Při formulování cílů ti pomůžou slovesa na straně \pageref{bloom}.}
\vspace*{1cm}

Jak poznám, že se cíle podařilo dosáhnout? \\
\note{Stanov si cíle a pak veď výuku tak, abys je naplnil(a).}

\hspace*{-1cm}
\rule{\rulelength}{0.4pt}

\marginsection{90}{-1.2}{po výuce}

\smileys

\selfassessment

\note{Sem napiš otázku, na kterou si chceš po výuce odpovědět.}\\
Vlastní otázka:

\newpage
\marginsection{-90}{1}{volné poznámky}
\vspace*{-2em}
\note{Podařilo se mi naplnit cíl hodiny?\\
Co je třeba do budoucna na téhle hodině změnit?\\
Jak budu plánovat a strukturovat nejbližší hodinu?\\
(Další tipy na otázky najdeš na straně \pageref{indikatory}.)}

% ======= WEEK 6 =======

\chapter{}
\vspace*{-2em}

\marginsection{90}{-1}{před výukou}

Příprava:\hspace{1cm} hodin(y)

Jaký je cíl této hodiny? Proč chci studenty učit tohle? \\
\note{Při formulování cílů ti pomůžou slovesa na straně \pageref{bloom}.}
\vspace*{1cm}

Jak poznám, že se cíle podařilo dosáhnout? \\
\note{Stanov si cíle a pak veď výuku tak, abys je naplnil(a).}

\hspace*{-1cm}
\rule{\rulelength}{0.4pt}

\marginsection{90}{-1.2}{po výuce}

\smileys

\selfassessment

\note{Sem napiš otázku, na kterou si chceš po výuce odpovědět.}\\
Vlastní otázka:

\newpage
\marginsection{-90}{1}{volné poznámky}
\vspace*{-2em}
\note{Podařilo se mi naplnit cíl hodiny?\\
Co je třeba do budoucna na téhle hodině změnit?\\
Co je problematické na úkolech, které jsem zadal?\\
(Další tipy na otázky najdeš na straně \pageref{indikatory}.)}

% ======= WEEK 7 =======
\commonweek

% ======= WEEK 8 =======
\commonweek

% ======= WEEK 9 =======
\commonweek

% ======= WEEK 10 =======
\commonweek

% ======= WEEK 11 =======
\commonweek

% ======= WEEK 12 =======

\chapter{}
\vspace*{-2em}

\marginsection{90}{-1}{před výukou}

Příprava:\hspace{1cm} hodin(y)

Jaký je cíl této hodiny? Proč chci studenty učit tohle?
\vspace*{1cm}

Jak poznám, že se cíle podařilo dosáhnout?

\hspace*{-1cm}
\rule{\rulelength}{0.4pt}

\marginsection{90}{-1.2}{po výuce}

\smileys

\selfassessment

Jaký je širší kontext toho, co učím?\\
\note{Kde v mojí výuce se dotýkám jiné výuky, jiných oblastí vědění?}

\newpage
\marginsection{-90}{1}{volné poznámky}
\vspace*{-2em}
\note{Jak získám od studentů zpětnou vazbu na ty aspekty mé vyúky, které mě zajímají?\\
Jak to dělají ostatní vyučující?\\
A jak jiné předměty?}

% ======= WEEK 13 =======

\chapter{}
\vspace*{-2em}

\marginsection{90}{-1}{před výukou}

Příprava:\hspace{1cm} hodin(y)

Jaký je cíl této hodiny? Proč chci studenty učit tohle?
\vspace*{1cm}

Jak poznám, že se cíle podařilo dosáhnout?

\hspace*{-1cm}
\rule{\rulelength}{0.4pt}

\marginsection{90}{-1.2}{po výuce}

\smileys

\selfassessment

Co mi dal tento semestr výuky?\\
\note{Co jsem si uvědomil(a)? Co jsem zlepšil(a)?}

\newpage
\marginsection{-90}{1}{volné poznámky}
\vspace*{-2em}
\note{Za co jsou studenti v předmětu hodnocení?\\
Odpovídá to, co zkouška testuje tomu, co já učím?\\
Pokud mají studenti zájem se dál samostatně rozvíjet v dané oblasti, jsou na to po tomto kurzu připraveni?}

% ======= APPENDICES =======

\restoregeometry
\chapter*{Rubrika}

TODO Uvodny text k rubrike -- co to je, kedy a ako sa to ma pouzit.
\begin{itemize}
\item Co to je?
\item Preco su tam 3 kategorie a co znamenaju?
\item Ako to pouzit? (na zaciatku a na konci roka si vyplnit vsetky skaly, kde sa aktualne nechadzam? Vybrat si 1-2 oblasti, kde sa chcem zlepsit, vymysliet konkretne kroky, najst si indikatory na sledovanie stavu, napisat si ich do stran k tyzdnom?)
\item Ine pouzitie: Moznost precitat si Mastery popisy a uvedomit si, co vsetko ucitel robi/mal by robit.
\item Dolezite je zamysliet sa nad sebou a svojim ucenim, ked to vyplnam!
\end{itemize}

\newpage
\rubricpage{Vědomá pozornost a rozhodování, \textit{tracking}}
{Nedokážu si při výuce ani po ní vybavit, co se děje a jaký mám záměr. Při výuce se často cítím ztracený.}
{Někdy se mi daří uvědomovat si, jaký mám právě záměr, co právě dělám a jaký to bude mít efekt. Většinou ale nemám při výuce vědomou pozornost.}
{Během výuky jsem si téměř vždy vědom toho, co se se mnou i se skupinou děje, co dělám a jaký to bude mít efekt, kam směřuji. Jsem schopen si vybavit, jak jsme se dostali do tam, kde právě jsme.}
{Interakce se skupinou, kladení otázek skupině}
{Se skupinou neinteraguji. Neptám se, asi bych nedostal odpověď. Nevím, jak studenty dobře zapojit.}
{Vím, že lze interagovat se skupinou a znám nástroje (chápu, jak by to šlo). Nejsem ale schopen je často používat. Někdy  se skupiny ptám, ale nedostávám odpověď.}
{Se skupinou často a efektivně interaguji, ptám se způsobem, který studenty aktivizuje a zapojuje. Situace, kdy nedostávám ze skupiny odpověď dokážu obratně řešit (např. přeformulováním otázky).}

\rubricpage{Strukturování výuky}
{O strukturování výuky nijak nepřemýšlím.}
{Chápu smysl přehledného strukturování své hodiny a snažím se o to. Často se ale zamotám, ztratím nit, nebo řeším mnoho naráz a studenti se pak ztrácí nebo odpojují.}
{Moje hodiny mají jasnou strukturu. Studenti ví co se právě děje, co bude následovat a chápou návaznosti. Mezi jednotlivými bloky vědomě dělám zřetelné přechody.}
{Vlastní pocit z výuky}
{Svou výuku nijak nereflektuji.}
{Ve výuce si často moc nevěřím, vyžaduje to hodně energie, nebo se cítím napjatý. Mám strach z toho, na co se studenti zeptají.}
{Ve výuce se většinou cítím uvolněně a sebevědomě, mám svůj styl a baví mě to.}

\rubricpage{Formativní zpětná vazba}
{O dávání zpětné vazby studentům nijak nepřemýšlím.}
{Snažím se studentům zpětnou vazbu dávat. Myslím si však, že jí není dost, nebo to nedělám efektivně, nebo to studenti nevnímají jako podporu a projev respektu.}
{Se studenty ve výuce interaguji tak, že dostávají průběžně formativní zpětnou vazbu. Chápou tedy, co jim jde, kde dělají chyby a jak se mohou zlepšovat. Zároveň se moji studenti cítí být respektováni a zpětné vazby se nebojí.}
{Jasné zadávání úloh}
{O zadávání úloh nijak zvlášť nepřemýšlím.}
{Stává se mi, že zadám úlohu a studenti neví, co dělat, jak začít nebo k čemu mají dojít (co má být výsledkem).}
{Když zadávám úlohu nebo aktivitu, mají studenti jasno v tom, co mají dělat a neřeší tak zbytečně věci, které nemám záměr procvičovat.}

\rubricpage{Design výuky a variabilita aktivit}
{Učím tak, jak mi řekli nebo kopíruji výuku, kterou jsem viděl. Nepřemýšlím o~jiných variantách.}
{Jsem si vědom toho, že existuje mnoho typů aktivit, které lze při výuce použít. Neznám jich ale dostatek, nejsem je schopen efektivně zadávat, nebo nemám jasno v tom proč je vybírat.}
{Znám dostatek typů aktivit, svou výuku skládám tak, aby byla dostatečně pestrá. Vybrané aktivity efektivně učí/procvičují to, co mám záměr procvičovat. Vybrané aktivity zároveň studenty efektivně zapojují a přispívají k jejich motivaci.}
{Širší kontext výuky a konkrétní hodiny}
{O širším kontextu výuky nepřemýšlím.}
{Uvědomuji si, že nemám pro sebe pojmenované znalosti a dovednosti, které u~studentů rozvíjím. Nevím, jak jejich pokrok sledovat. Nevím, v jakém kontextu to studenti využijí.}
{Mám jasnou představu o tom, k čemu studenty vedu, jaké dovednosti rozvíjím, jaké znalosti jim chci předat. Vím, proč tyto dovednosti rozvíjím a v jakém kontextu je studenti v budoucnu použijí. Vím, jak jejich pokrok sledovat.}
\vspace*{-1em}

\rubricpage{Jasné vysvětlování}
{Svoje vysvětlování nijak nereflektuji.}
{Když vysvětluji, běžně se mi stává, že si nejsem jistý, zda vysvětluji dobře a zda to studentům pomáhá při pochopení.}
{Když vysvětluji teorii, demonstruji řešení nebo opravuji chybný postup. Dělám to efektivně a dokážu se dobře vžít do toho, jak to vidí student. Nestává se mi, že by studenti moje vysvětlení nechápali, nebo že bych vysvětloval něco, na co se neptali.}
{Nastavení prostředí, systémy ve výuce}
{O nastavení atmosféry nepřemýšlím, systémy ve své výuce nevnímám.}
{Přemýšlím o nastavení pravidel i atmosféry. Systémy ve výuce (např. bodování) přebírám od ostatních. Nemám ale jasno v efektech, nebo neznám způsoby, jak bych je mohl upravit.}
{Umím ve výuce vytvořit prostředí, které podporuje efektivní učení. U systémů, které používám (např. bodování, bonbóny, zahajovací rituály) chápu efekt. Systémy nepřebírám slepě, chápu jejich účel a uzpůsobuji podle svých potřeb.}

\rubricpage{Improvizace, přizpůsobování}
{V průběhu své výuky na vzniklé situace nijak vědomě nereaguji.}
{Uvědomuji si chvíle, kdy by mohlo být zajímavé nebo užitečné dělat něco jiného, než jsem měl v plánu. Většinou ale nedokážu v daný moment vhodně zareagovat.}
{Jsem schopen svoji výuku průběžně přizpůsobovat tomu, co se právě děje ve skupině a co studenti potřebují. Mám k tomu dostatek nástrojů a dokážu je efektivně použít.}
{Individuální interakce se studenty}
{O konzultacích nijak vědomě nepřemýšlím.}
{Stává se mi, že si nevím rady při individuální interakci se studentem (např. u~tabule, konzultace). Interakce neprobíhá efektivně nebo se student cítí zastrašen.}
{Při interakcích s jednotlivci (např. u tabule, konzultace) efektivně využívám čas. Studenti se mnou konzultují rádi, cítí ze mě respekt a podporu.}

\rubricpage{Management skupiny, \textit{subgrouping}}
{Přijde mi, že dělit skupinu je zbytečnost, o~práci v~menších skupinkách nepřemýšlím.}
{Uvědomuji si možnosti práce ve dvojicích či menších skupinách. Tuším, že bych toho mohl lépe využívat, hledám jak.}
{Mám jasno v tom, kdy chci pracovat se skupinou jako celkem, kdy s jednotlivci a kdy se skupinkami. Rozdělení do menších skupin efektivně používám. Ve vhodných případech zadávám interakce mezi skupinami.}
{Čtení skupiny}
{Skupinu ve své výuce nesleduji, pozornost věnuji pouze obsahu.}
{Uvědomuji si, že mi skupina vysílá signály a že by bylo dobré jim rozumět a využít je pro efektivní vedení výuky. Ve výuce to ale dokážu jen výjimečně.}
{Dokážu dobře odhadnout naladění skupiny. Mám jasno v tom, co se ve skupině studentů děje (např. únava, rezignace, zájem).}

\chapter*{Indikátory}
\label{indikatory}
\vspace{-0.5em}
\note{Ukazatel stavu; aspekt výuky, který je možné přímo sledovat.}

\section*{Kvantitativní indikátory}

Práce s časem\\
\note{O kolik minut jsem prodlužoval(a)? Kolik času zůstalo navíc?\\
Jak dlouho jsem měl(a) monolog? Kolik času studenti něco dělali?\\
Kolik času jsem v průměru věnoval(a) jednomu studentovi?}

Interakce se studenty\\
\note{Kolik otázek jsem položil(a)? Kolik z nich bylo uzavřených/otevřených?\\
Kolik otázek jsem dostal(a)? Kolik jsem dokázal(a) zodpovědět?\\
Kolik studentů jsem pochválil(a) (za jakkoliv malý úspěch)?}

Řešení úloh\\
\note{Kolik úloh vyřešili studenti na hodině? Kolik úkolů dostali na doma?\\
Kolik studentů bylo ztraceno (těžký úkol)? Kolik se naopak nudilo?}

Účast\\
\note{Kolik studentů přišlo na výuku? Kolik z nich meškalo?\\
Kolik studentů odešlo během výuky?}

\section*{Kvalitativní indikátory}

Pocity a subjektivní dojmy\\
\note{Jaké pocity jsem měl(a) během výuky?\\
Jací byli studenti? Byli aktivní, poslouchali mě, \dots?\\
Dávám svojí prací studentům dobrý příklad?}

Organizace\\
\note{Věděli studenti, co přesně měli dělat, jak, a proč?\\
Do jaké míry jsem dodržel(a) časový plán?}

Obsah výuky\\
\note{Je moje výuka variabilní (různé typy úloh, pomůcek, \dots)?\\
Odpovídá obsah mojí výuky tomu, co chci naučit?}

\newpage

\chapter*{Bloomova taxonomie}
\label{bloom}
\vspace{-0.5em}
\note{Hierarchie kognitivních vzdělávacích cílů, B.\ Bloom, 1956}
\vspace{-0.3em}

\begin{enumerate}[leftmargin=*]
\item \textbf{Vzpomenout si (\textit{remember})}\\
\note{termíny a fakta, jejich klasifikace a kategorizace}\\
{\small definovat, vyjmenovat, zopakovat, popsat, rozpoznat, přiřadit pojem ke konceptu}

\item \textbf{Pochopit (\textit{understand})}\\
\note{překlad do jiné formy, jednoduchá interpretace, extrapolace}\\
{\small interpretovat, reprezentovat, vysvětlit, parafrázovat, sumarizovat, uvést příklad, ilustrovat na příkladě, porovnat}

\item \textbf{Použít (\textit{apply})}\\
\note{užití postupu ve správné situaci, použití abstrakcí a zobecnění}\\
{\small vykonat, aplikovat, využít, demonstrovat v praxi, implementovat, vyřešit uzavřený a ohraničený problém}

\item \textbf{Analyzovat (\textit{analyze})}\\
\note{rozklad na části, vztahy a interakce mezi částmi}\\
{\small diskutovat, \uv{rozbít} na části, rozlišit, navrhnout/zvolit řešení, organizovat, strukturovat, integrovat, propojit}

\item \textbf{Vyhodnotit (\textit{evaluate})}\\
\note{posouzení na základě stanovených kritérií a standardů}\\
{\small zkontrolovat, detekovat, testovat, posoudit, odborně kritizovat, argumentovat}

\item \textbf{Vytvořit (\textit{create})}\\
\note{vytvoření nového celku, reorganizace do nové struktury}\\
{\small generovat, produkovat, vymyslet, plánovat/projektovat, komponovat, designovat, budovat, kreativně kombinovat prvky}
\end{enumerate}

\note{Hranice mezi úrovněmi nejsou striktně určené. Seznam akcí ti ale pomůže popsat cíle hodiny i znalosti a dovednosti, které by měli studenti získat. Ukáže ti také, jestli je tvá výuka dostatečně variabilní.}

\chapter*{Užitečné nástroje}

\note{Seznam výukových nástrojů, které se můžou hodit.}

\section*{Základní nástroje}
\begin{itemize}
\item Otázky -- \note{srozumitelnost, požadovaná akce, poděkování}
\item Samostatná práce -- \note{první krok, požadovaný výsledek, odhad trvání}
\item Domácí úkoly -- \note{způsob odevzdání, hodnocení, cíl}
\end{itemize}

\section*{Alternativní nástroje dotazování se}
\begin{itemize}
\item Hlasování -- \note{pozitivní, negativní, rozdělování bodů}
\item Kahoot -- \note{rychlý celkový přehled, anonymita, cíl}
\end{itemize}

\section*{Alternativní nástroje samostatné práce}
\begin{itemize}
\item Think-Pair-Share -- \note{nejprve sami, diskuze v páru, sklízení}
\item Anti-problém -- \note{řešení opaku problému}
\item Myšlenkový oblak -- \note{sbírání asociací, diskuze, souvislosti}
\item Rozcvička -- \note{fyzické/myšlenkové probuzení, opakování}
\item Figurky -- \note{manipulace s fyzickými předměty}
\end{itemize}

\section*{Alternativní nástroje domácích úkolů}
\begin{itemize}
\item Externí motivace -- \note{body navíc, bonbóny, ...}
\item Peer-review -- \note{učení se učením/hodnocením}
\end{itemize}

\section*{Základní nástroje vysvětlování}
\begin{itemize}
\item Časový plán -- \note{odhad trvání}
\item Interaktivita -- \note{otázky, diskuze, vzájemné vysvětlování}
\item Zpětné vazby -- \note{ujištění se, že studenti pochopili správně to, co se řeklo}
\item Praktická ukázka -- \note{aby se museli i dívat, nejen poslouchat}
\item Slajdy (případně checkboxy na tabuli)\\ -- \note{vizualizace aktuálního pokroku hodiny}
\item Precedenty -- \note{atmosféra, dochvilnost, úkoly, otázky}
\item Bloomova taxonomie -- \note{verbalizace cílu hodiny/výuky}
\item Veřejný checklist -- \note{přehledná struktura, orientace}
\item Schody -- \note{postupnost ztěžujících se úkolů}
\item Rubrika -- \note{uvědomění schopností, sledování pokroku}
\item Nastavování prahu -- \note{Jak moc snadné je to udělat?}
\item Alternativní zadání -- \note{programování na 4 způsoby}
\item Hospitace -- \note{Jak to dělají jiní? Jak to dělám já?}
\item Argumentace řešení -- \note{...}
\item Konceptová mapa -- \note{vizualizace struktury znalostní báze}
\item Hledání rozdílů -- \note{...}
\item Subgrouping -- \note{...}
\item Check-in -- \note{...}
\item Check-out -- \note{...}
\item Tracking -- \note{...}
\item Uspořádání prostoru -- \note{rovnost, tvar usazení, pozice učitele}
\end{itemize}

\chapter*{Prostor pro poznámky}
\note{Například nejdůležitější feedback z hospitací.}

\chapter*{Prostor pro poznámky}
\note{Na co chci myslet, když učím?}

\newpage
\vspace*{\fill}
Verze: \today

Vytvořeno pro kurz \textit{DUCIT Praktikum vedení cvičení} na Fakultě informatiky MU.

Autoři deníku: Martin Ukrop, Ondráš Přibyla, Martin Macák, Valdemar Švábenský

\textit{Budeme rádi, když nám dáš vědět svoje postřehy z používaní deníku a další náměty emailem Martinovi na mukrop@mail.muni.cz.}
\end{document}
